\documentclass[12pt]{article}
\usepackage{amsmath, amssymb, amsthm}

\newtheorem{definition}{Definition}
\newtheorem{proposition}{Proposition}
\newtheorem{theorem}{Theorem}
\newtheorem{example}{Example}

\begin{document}
	
	\section*{Cyclic Groups: Introduction}
	
	A \emph{cyclic group} is a group generated by a single element $g$.  
	This means every element of the group can be written as a power or multiple of $g$:
	
	\[
	\langle g \rangle =
	\begin{cases}
		\{\, g^n \mid n \in \mathbb{Z} \,\}, & \text{(multiplicative notation)} \\[6pt]
		\{\, n g \mid n \in \mathbb{Z} \,\}, & \text{(additive notation)} 
	\end{cases}
	\]
	
	- **Identity element**:  
	$g^0 = e$ in multiplicative notation, and $0g = e$ in additive notation.  
	For the group of complex roots of unity, the identity is $1 \in \mathbb{C}$.
	
	- **Infinite vs finite**:  
	\[
	\langle g \rangle \cong (\mathbb{Z}, +) \quad \text{(infinite case)}, 
	\]
	\[
	\langle g \rangle \cong (\mathbb{Z}_n, +) \quad \text{(finite case, $g$ of order $n$)}.
	\]
	
	- **Subgroups**:  
	Every subgroup of a cyclic group is cyclic.  
	If $|G| = n$, then each divisor $d \mid n$ determines a unique subgroup of order $d$.
	
	- **Structural role**:  
	Cyclic groups are the basic building blocks of abelian groups:  
	every finite abelian group can be expressed as a product of cyclic groups.
	
	\subsection*{Example: The Cyclic Group $\mathbb{Z}_6$}
	
	Consider $(\mathbb{Z}_6, +)$, the integers modulo $6$ under addition.
	
	\[
	\langle 1 \rangle = \{0,1,2,3,4,5\} = \mathbb{Z}_6,
	\]
	so $1$ is a generator of the entire group.
	
	Other elements may generate proper subgroups:
	\[
	\langle 2 \rangle = \{0,2,4\}, \qquad 
	\langle 3 \rangle = \{0,3\}.
	\]
	
	Thus, the subgroups of $\mathbb{Z}_6$ correspond exactly to the divisors of $6$.
	
	
	\section*{Cyclic Groups}
	
	\subsection*{Group Axioms}
	
	A group is a pair $(G,*)$ where $G$ is a set and $*$ is a binary operation on $G$ such that:
	
	\begin{enumerate}
		\item \textbf{Closure:} For all $a,b \in G$, $a*b \in G$.
		\item \textbf{Associativity:} For all $a,b,c \in G$, $(a*b)*c = a*(b*c)$.
		\item \textbf{Identity:} There exists $e \in G$ such that $e*a = a*e = a$ for all $a \in G$.
		\item \textbf{Inverses:} For each $a \in G$, there exists $a^{-1} \in G$ such that $a*a^{-1} = a^{-1}*a = e$.
	\end{enumerate}
	
	\subsection*{Cyclic Groups}
	
	\begin{definition}
		A group $G$ is \emph{cyclic} if there exists an element $g \in G$ such that
		$$
		G = \langle g \rangle = \{ g^n : n \in \mathbb{Z} \}
		$$
		in multiplicative notation, or
		$$
		G = \langle g \rangle = \{ n g : n \in \mathbb{Z} \}
		$$
		in additive notation. Such an element $g$ is called a \emph{generator} of $G$.
	\end{definition}
	
	\subsection*{Examples}
	
	\begin{example}[Infinite Cyclic Group]
		In $(\mathbb{Z},+)$, the element $1$ generates all integers:
		$$
		\langle 1 \rangle = \{ n \cdot 1 : n \in \mathbb{Z} \} = \mathbb{Z}.
		$$
		Thus $(\mathbb{Z},+)$ is an infinite cyclic group, with generators $\pm 1$.
	\end{example}
	
	\begin{example}[Finite Cyclic Group]
		In $(\mathbb{Z}_n,+)$, an element $a$ is a generator if and only if $\gcd(a,n)=1$. 
		The number of generators is $\varphi(n)$, Euler's totient function.
	\end{example}
	
	\begin{example}[Roots of Unity]
		The group of $n$-th roots of unity is
		$$
		\mu_n = \{ e^{2\pi i k / n} : k = 0,1,\dots,n-1 \} \subset \mathbb{C}^\times.
		$$
		This is a cyclic group of order $n$. A primitive $n$-th root of unity is an element $\zeta = e^{2\pi i k / n}$ with $\gcd(k,n)=1$. 
		The primitive roots are exactly the generators of $\mu_n$, and there are $\varphi(n)$ of them.
	\end{example}
	
	\subsection*{Classification of Cyclic Groups}
	
	\begin{theorem}
		Every cyclic group is isomorphic to either $(\mathbb{Z},+)$ (infinite case) or $(\mathbb{Z}_n,+)$ (finite case).
	\end{theorem}
	
	\subsection*{Subgroups of Cyclic Groups}
	
	\begin{theorem}
		Every subgroup of a cyclic group is cyclic.
	\end{theorem}
	
	\begin{proposition}
		If $G = \langle g \rangle$ is cyclic of order $n$, then for each divisor $d \mid n$ there exists a unique subgroup of order $d$, namely
		$$
		\langle g^{n/d} \rangle.
		$$
	\end{proposition}
	
	\subsection*{Order of Elements}
	
	\begin{theorem}
		If $G = \langle g \rangle$ has order $n$, then
		$$
		\operatorname{ord}(g^k) = \frac{n}{\gcd(n,k)}.
		$$
		In particular, $g^k$ is a generator if and only if $\gcd(n,k)=1$. 
		Thus the number of generators of $G$ is $\varphi(n)$.
	\end{theorem}
	
	\begin{example}[Primitive vs. non-primitive roots for $n=8$]
		
		The $8$-th roots of unity are
		$$
		\mu_8 = \{ e^{2\pi i k / 8} : k=0,1,\dots,7 \}
		= \{ 1, \, e^{2\pi i/8}, \, e^{2\pi i\cdot 2/8}, \dots, e^{2\pi i\cdot 7/8} \}.
		$$
		
		\begin{itemize}
			\item Take $\zeta = e^{2\pi i/8}$. Since $\gcd(1,8)=1$, $\zeta$ is a primitive $8$-th root of unity.
			Its powers generate all eight roots:
			$$
			\langle \zeta \rangle = \{ 1, \zeta, \zeta^2, \zeta^3, \zeta^4, \zeta^5, \zeta^6, \zeta^7 \} = \mu_8.
			$$
			
			\item Take $\eta = e^{2\pi i \cdot 2/8} = e^{\pi i/2} = i$. 
			Here $\gcd(2,8)=2 \neq 1$, so $\eta$ is not primitive. 
			Its powers generate only the $4$-th roots of unity:
			$$
			\langle \eta \rangle = \{ 1, i, -1, -i \} \subsetneq \mu_8.
			$$
		\end{itemize}
		
		Thus, a primitive root generates all $n$-th roots of unity, while a non-primitive root generates only a proper subgroup corresponding to a divisor of $n$.
	\end{example}
	\section*{Roots of Unity}
	
	\begin{definition}
		A complex number $z$ is called a \emph{root of unity} if there exists $n \in \mathbb{N}$ such that
		$$
		z^n = 1.
		$$
		The set of all $n$-th roots of unity is
		$$
		\mu_n = \{ e^{2\pi i k / n} : k = 0,1,\dots,n-1 \}.
		$$
	\end{definition}
	
	\subsection*{Examples}
	
	\begin{example}[Square roots of unity]
		The equation $z^2=1$ has solutions
		$$
		\mu_2 = \{1, -1\}.
		$$
	\end{example}
	
	\begin{example}[Cube roots of unity]
		The equation $z^3=1$ has solutions
		$$
		\mu_3 = \{ 1, \ e^{2\pi i/3}, \ e^{4\pi i/3} \}.
		$$
		In explicit form,
		$$
		\mu_3 = \left\{ 1, \ -\tfrac{1}{2}+\tfrac{\sqrt{3}}{2}i, \ -\tfrac{1}{2}-\tfrac{\sqrt{3}}{2}i \right\}.
		$$
		These are the vertices of an equilateral triangle on the unit circle.
	\end{example}
	
	\begin{example}[Fourth roots of unity]
		The equation $z^4=1$ has solutions
		$$
		\mu_4 = \{ 1, i, -1, -i \}.
		$$
		These are the four vertices of a square on the unit circle.
	\end{example}
	
	\subsection*{Primitive Roots of Unity}
	
	A root of unity $\zeta \in \mu_n$ is \emph{primitive} if its order is exactly $n$, i.e.
	$$
	\langle \zeta \rangle = \{1, \zeta, \zeta^2, \dots, \zeta^{n-1}\} = \mu_n.
	$$
	
	\begin{example}[Primitive vs. non-primitive]
		In $\mu_4$:
		\begin{itemize}
			\item $\zeta = i = e^{2\pi i/4}$ is primitive, since
			$$
			\langle i \rangle = \{1,i,-1,-i\} = \mu_4.
			$$
			\item $\eta = -1$ is not primitive, since
			$$
			\langle -1 \rangle = \{1,-1\} \subsetneq \mu_4.
			$$
		\end{itemize}
	\end{example}
	
	\section*{Worked Example: The Cyclic Group $(\mathbb{Z}_6,+)$}
	
	\subsection*{Basic Data}
	Elements:
	$$
	\mathbb{Z}_6=\{\,0,1,2,3,4,5\,\}.
	$$
	Additive order formula:
	$$
	\operatorname{ord}(\overline{a})=\frac{6}{\gcd(a,6)}.
	$$
	
	\subsection*{Element–Order–Generator Summary}
	\begin{center}
		\begin{tabular}{c|c|l|c}
			$a$ & $\operatorname{ord}(\overline{a})$ & $\langle \overline{a}\rangle$ & generator? \\ \hline
			0 & 1 & $\{\,0\,\}$ & no \\
			1 & 6 & $\{\,0,1,2,3,4,5\,\}=\mathbb{Z}_6$ & yes \\
			2 & 3 & $\{\,0,2,4\,\}$ & no \\
			3 & 2 & $\{\,0,3\,\}$ & no \\
			4 & 3 & $\{\,0,4,2\,\}=\{\,0,2,4\,\}$ & no \\
			5 & 6 & $\{\,0,5,4,3,2,1\,\}=\mathbb{Z}_6$ & yes \\
		\end{tabular}
	\end{center}
	
	Generators are exactly those $a$ with $\gcd(a,6)=1$:
	$$
	\{\text{generators of }(\mathbb{Z}_6,+)\}=\{\,1,5\,\}.
	$$
	
	\subsection*{Canonical Subgroups (one for each divisor of $6$)}
	Divisors of $6$ are $1,2,3,6$. The corresponding unique subgroups are:
	\begin{align*}
		&d=1:\quad \text{order }1,\ \langle \overline{0}\rangle=\{\,0\,\}.\\[4pt]
		&d=2:\quad \text{order }2,\ \langle \overline{3}\rangle=\{\,0,3\,\}.\\[4pt]
		&d=3:\quad \text{order }3,\ \langle \overline{2}\rangle=\langle \overline{4}\rangle=\{\,0,2,4\,\}.\\[4pt]
		&d=6:\quad \text{order }6,\ \langle \overline{1}\rangle=\langle \overline{5}\rangle=\mathbb{Z}_6.
	\end{align*}
	
	\subsection*{Cayley Table of $(\mathbb{Z}_6,+)$}
	\begin{center}
		\begin{tabular}{c|cccccc}
			$+$ & 0 & 1 & 2 & 3 & 4 & 5 \\ \hline
			0 & 0 & 1 & 2 & 3 & 4 & 5 \\
			1 & 1 & 2 & 3 & 4 & 5 & 0 \\
			2 & 2 & 3 & 4 & 5 & 0 & 1 \\
			3 & 3 & 4 & 5 & 0 & 1 & 2 \\
			4 & 4 & 5 & 0 & 1 & 2 & 3 \\
			5 & 5 & 0 & 1 & 2 & 3 & 4 \\
		\end{tabular}
	\end{center}
	
	\subsection*{Checks}
	\begin{itemize}
		\item Identity: $0$; inverse of $a$ is $-a\equiv 6-a\pmod{6}$.
		\item Orders match the formula:
		$$
		\operatorname{ord}(\overline{0})=1,\ 
		\operatorname{ord}(\overline{1})=\operatorname{ord}(\overline{5})=6,\ 
		\operatorname{ord}(\overline{2})=\operatorname{ord}(\overline{4})=3,\ 
		\operatorname{ord}(\overline{3})=2.
		$$
		\item Subgroups correspond to divisors of $6$ and are unique for each divisor (cyclic subgroup structure).
	\end{itemize}
	
	
\end{document}
